% =============================================================================
% AI and Machine Learning -- Project Proposal Template
% Alpen-Adria-Universität Klagenfurt
%
% Page limit: 4 pages, EXCLUDING references. Everything else counts,
% including figures, tables, the appendix and the AI disclosure statement.
% Do not change margins, font sizes or spacing.
%
% The comments under each heading describe what the section must contain.
% Replace every "Your text." with your content and submit the PDF in Moodle.
% The five sections below are required, in this order.
%
% Projects are done in teams of 3 students. The team shares one dataset and
% one learning task (T, P, E); each student formulates, trains, tunes and
% evaluates a fundamentally distinct model family. State in every section who
% is expected to do what, e.g. "(name) will implement the logistic regression
% baseline, ...".
% =============================================================================

\documentclass{article}

% "preprint" shows your names; without it the template produces an
% anonymous submission.
\usepackage[preprint]{neurips_2025}

\usepackage[utf8]{inputenc} % allow utf-8 input
\usepackage[T1]{fontenc}    % use 8-bit T1 fonts
\usepackage{hyperref}       % hyperlinks
\usepackage{url}            % simple URL typesetting
\usepackage{booktabs}       % professional-quality tables
\usepackage{amsmath}        % mathematical environments
\usepackage{amsfonts}       % blackboard math symbols
\usepackage{nicefrac}       % compact symbols for 1/2, etc.
\usepackage{microtype}      % microtypography
\usepackage{xcolor}         % colors
\usepackage{graphicx}       % figures

\title{Project Proposal: Your Project Title}

\author{%
  Student A \\
  Student ID: 00000000 \\
  \texttt{student.a@edu.aau.at} \\
  \And
  Student B \\
  Student ID: 00000000 \\
  \texttt{student.b@edu.aau.at} \\
  \And
  Student C \\
  Student ID: 00000000 \\
  \texttt{student.c@edu.aau.at} \\
}

\begin{document}

\maketitle

% =============================================================================
\section{Motivation and Domain Problem}
% - Real-world context and personal motivation.
% - Domain significance: which decisions or actions depend on these predictions?
% - Hypothesis: why is this learning task solvable from empirical data?

Your text.

% =============================================================================
\section{Learning Task $(T, P, E)$}

\subsection{Task $T$}
% - Formal type of the task (binary classification, multi-class classification,
%   regression, clustering) and definition of the target variable.

Your text.

\subsection{Performance Measure $P$}
% - The evaluation measures you choose, justified by the class balance and the
%   domain costs of false positives vs. false negatives.

Your text.

\subsection{Experience $E$}
% - Data source, sample size N, number of features D, variable types
%   (continuous, ordinal, categorical), and representative data samples.

Your text.

% =============================================================================
\section{Previous Work and Literature Baselines}
% - Summarize at least 3 related publications, benchmarks or authoritative
%   studies that tackle similar tasks.
% - Analyze their methods and reported performance, and explain how their
%   findings serve as a baseline for your project.
% - Cite your sources, e.g. \citet{james2023islp} or \citep{bishop2024deep}, and
%   add them to references.bib.

Your text.

% =============================================================================
\section{Implementation Plan and Individual Roles}
% - Planned ML pipeline: data cleaning, feature encoding, scaling, model
%   training and cross-validation strategy.
% - Distribution of the model families across the 3 students, e.g. in a table
%   like Table~\ref{tab:roles}. Each student takes a fundamentally distinct
%   model family.

Your text.

\begin{table}[h]
  \caption{Model families of the team members.}
  \label{tab:roles}
  \centering
  \begin{tabular}{ll}
    \toprule
    Team member & Model family \\
    \midrule
    Student A & e.g.\ linear or regularized model (logistic regression, ElasticNet) \\
    Student B & e.g.\ tree ensemble (random forest, gradient boosting) \\
    Student C & e.g.\ support vector machine (RBF kernel) or MLP \\
    \bottomrule
  \end{tabular}
\end{table}

% =============================================================================
\section{Risk Management}
% - Anticipated bottlenecks, e.g. data scarcity, extreme class imbalance,
%   non-converging optimization, computational runtime limits.
% - Concrete fallback strategies ("plan B") if your primary hypotheses fail.

Your text.

% =============================================================================
\bibliographystyle{plainnat}
\bibliography{references}

% =============================================================================
\appendix

\section{AI Disclosure Statement}
% - Declare any use of AI tools, e.g. for brainstorming, exploratory data
%   analysis, code generation, text drafting or literature search.
% - Specify the tools, model versions and the general purpose of their use.
% - If you did not use any, say so.

Your text.

\end{document}
